\section{Converg\^encia em $h$: o protocolo}
\label{sec:convergencia}

\subsection{Por que este estudo, e por que agora}

A Se\c{c}\~ao~\ref{sec:vies} deixou uma pergunta com nome: os cinco desvios contra
os valores publicados s\~ao \emph{todos negativos}, em duas discretiza\c{c}\~oes
independentes. Isso \'e assinatura de vi\'es compartilhado, e tr\^es candidatos
foram nomeados sem que nenhum fosse medido --- a resolu\c{c}\~ao \'unica, a
restri\c{c}\~ao capilar de passo, e o refino da esteira poder ser insuficiente.

Um estudo de converg\^encia separa os tr\^es. Se os desvios encolhem com $h$ na
taxa esperada, o vi\'es \'e erro de discretiza\c{c}\~ao e o par de
implementa\c{c}\~oes est\'a correto; se n\~ao encolhem, a causa \'e outra, e o
alvo da investiga\c{c}\~ao muda de lugar.

Este relat\'orio listava converg\^encia em $h$ como pend\^encia desde a primeira
vers\~ao, por completude metodol\'ogica. Ela passa a ser a \emph{pergunta aberta},
por raz\~ao medida.

\subsection{O que o custo obriga: engrossar, e n\~ao refinar}

A escolha natural seria refinar. Ela \'e invi\'avel, e o motivo \'e a pr\'opria
f\'isica do problema: a restri\c{c}\~ao capilar de passo, $\Delta t <
\sqrt{\rho h^{3}/2\pi\sigma}$, escala como $h^{3/2}$. Dobrar a resolu\c{c}\~ao
multiplica as c\'elulas por 4 \emph{e} os passos por $2{,}8$:

\pgfplotstabletypeset[
  string type,
  columns/malha/.style={column name={malha}},
  columns/celraio/.style={column name={c\'el./raio}},
  columns/dt/.style={column name={$\Delta t$}},
  columns/passos/.style={column name={passos}},
  columns/celulas/.style={column name={c\'elulas}},
  columns/custo/.style={column name={custo}},
  columns/horas/.style={column name={horas}},
  every head row/.style={before row=\toprule, after row=\midrule},
  every last row/.style={after row=\bottomrule},
]{\dat{convergencia-custo.dat}}

\SI{136}{\hour} --- quase seis dias --- para \emph{um} m\'etodo, e o estudo pede
os dois. A malha refinada est\'a fora de alcance nesta m\'aquina.

Engrossar custa quase nada e serve igualmente para estimar ordem. O estudo usa
tr\^es pontos com raz\~ao de refino $\approx 1{,}41$ entre consecutivos --- quase
uniforme, o que \'e condi\c{c}\~ao para extrapola\c{c}\~ao de Richardson --- e
aproveita a corrida do fechamento como o ponto mais fino:

\pgfplotstabletypeset[
  string type,
  columns/corrida/.style={column name={corrida}},
  columns/celraio/.style={column name={c\'el./raio}},
  columns/dt/.style={column name={$\Delta t$}},
  columns/proposito/.style={column name={prop\'osito}},
  every head row/.style={before row=\toprule, after row=\midrule},
  every last row/.style={after row=\bottomrule},
]{\dat{convergencia-plano.dat}}

\subsection{O confundidor, e a corrida que existe para remov\^e-lo}

H\'a um problema de desenho que precisa ser dito antes dos resultados: como
$\Delta t$ \'e fixado pela restri\c{c}\~ao capilar, \textbf{engrossar a malha
tamb\'em engrossa o passo}. As corridas A, B e C variam os dois \emph{juntos}, e
portanto medem o efeito combinado --- n\~ao a resolu\c{c}\~ao isolada.

Sem cuidado, isso produziria uma conclus\~ao falsa de qualquer um dos dois lados:
um desvio que encolhesse ao refinar seria atribu\'ido \`a malha quando poderia ser
do passo, e um que n\~ao encolhesse absolveria a malha indevidamente.

A corrida D existe s\'o para isso: mesma malha de A, com $\Delta t$ pela metade.
A diferen\c{c}a entre A e D \'e efeito \emph{puro} de passo, na resolu\c{c}\~ao em
que ele \'e barato de medir. Se ela for desprez\'ivel diante da diferen\c{c}a
entre A e C, a varia\c{c}\~ao ao longo de A--B--C pode ser lida como efeito de
malha; se n\~ao for, as duas contribui\c{c}\~oes precisam ser separadas antes de
qualquer afirma\c{c}\~ao sobre ordem.

\subsection{As leituras, declaradas antes dos n\'umeros}
\label{sec:conv-leituras}

Como na Se\c{c}\~ao~\ref{sec:leituras}, a interpreta\c{c}\~ao fica fixada agora,
para que o resultado n\~ao seja explicado depois de conhecido. Os cinco desvios
ser\~ao medidos em A, B e C, e a leitura depende de \emph{como} eles variam:

\begin{itemize}
  \item \textbf{Encolhem monotonicamente com $h$, com ordem entre 1 e 2.} O vi\'es
    \'e erro de discretiza\c{c}\~ao, e as implementa\c{c}\~oes est\~ao corretas ---
    a compara\c{c}\~ao com a literatura passa a ser uma quest\~ao de
    resolu\c{c}\~ao, n\~ao de m\'etodo. \'E o desfecho esperado, e o menos
    informativo.
  \item \textbf{Encolhem, mas a extrapola\c{c}\~ao n\~ao chega \`a faixa
    publicada.} Sobra vi\'es residual que o refino n\~ao remove, e os outros dois
    candidatos --- contorno e refino da esteira --- passam a ser os suspeitos.
  \item \textbf{N\~ao encolhem, ou encolhem em ordem menor que 1.} A causa n\~ao \'e
    resolu\c{c}\~ao. \'E o desfecho mais informativo e o mais trabalhoso: aponta
    formula\c{c}\~ao ou condi\c{c}\~ao de contorno, e invalida a explica\c{c}\~ao
    mais confort\'avel.
  \item \textbf{Variam sem monotonicidade.} O estudo n\~ao decide nada e o
    problema \'e outro --- provavelmente a adapta\c{c}\~ao da malha mudando de
    padr\~ao entre resolu\c{c}\~oes, j\'a que a banda do crit\'erio \'e contada em
    \emph{c\'elulas} e portanto muda de largura f\'isica com $h$.
\end{itemize}

\subsection{O que este estudo n\~ao decide}

Convém delimitar antes, e n\~ao depois:

\begin{itemize}
  \item \textbf{S\~ao tr\^es pontos, e dois n\'iveis de AMR fixos.} A raz\~ao
    entre eles \'e $\approx 1{,}41$, e n\~ao $2$: a ordem estimada \'e mais
    sens\'ivel a ru\'ido do que seria com raz\~ao maior.
  \item \textbf{A banda do crit\'erio \'e contada em c\'elulas finas}, ent\~ao a
    regi\~ao refinada tem largura \emph{f\'isica} diferente em cada
    resolu\c{c}\~ao. O refino n\~ao \'e, a rigor, auto-semelhante --- e esse \'e
    exatamente o quarto desfecho previsto acima.
  \item \textbf{A resolu\c{c}\~ao mais fina \'e a atual}, n\~ao uma mais fina que
    ela. O estudo estima a ordem a partir de malhas \emph{mais grossas}, e
    extrapolar para $h\to0$ a partir dai' carrega mais incerteza do que
    extrapolar de dentro de uma s\'erie refinada.
  \item \textbf{Uma \'unica realiza\c{c}\~ao por ponto.} N\~ao h\'a barra de erro
    pr\'opria; a refer\^encia para julgar signific\^ancia continua sendo a
    dispers\~ao entre os tr\^es c\'odigos publicados.
\end{itemize}
